%! AP Statistics — Unit 4, Lesson 4.6 — Lesson Plan
\documentclass[10pt]{article}
\usepackage{apstats-article}
\usepackage{apstats-boxes}
\usepackage{graphicx}

\newcommand{\UnitNumberName}{Unit 4: Probability, Random Variables, and Probability Distributions \quad}
\newcommand{\LessonNumberName}{Lesson 4.6: Independent Events and Unions of Events}

\begin{document}
\begin{center}
    {\Large\bfseries \CourseName: \SchoolYear\ (\MeetingLength) \\
    {\small\bfseries \UnitNumberName \LessonNumberName}}
\end{center}

\begin{tcolorbox}[colback=sky, colframe=navy, boxrule=0.9pt, arc=2mm,
    left=2mm, right=2mm, top=1.5mm, bottom=1.5mm]
    \textbf{Primary Objective:} Students will determine whether two events are
    independent using conditional probabilities, and calculate the probability of
    the union of two events using the addition rule $P(A \cup B) = P(A) + P(B) - P(A \cap B)$
    (Big Idea: VAR-4).
\end{tcolorbox}

\renewcommand{\arraystretch}{1.6}
\begin{skillbox}[Priority Ideas \& Skills]{goldbox}
\begin{minipage}[t]{0.48\linewidth}
    \textbf{AP Skill 3.A --- Determine Parameters for Probability Distributions}
    \begin{itemize}
        \item Test for independence: $P(A|B) = P(A)$, or equivalently
              $P(A \cap B) = P(A) \cdot P(B)$ (VAR-4.E.1, VAR-4.E.2).
        \item Calculate $P(A \cup B)$ using the addition rule (VAR-4.E.4).
        \item Distinguish independence from mutually exclusive --- the most
              common AP exam confusion in this unit.
        \item Verify independence from a two-way table.
    \end{itemize}
\end{minipage}
\hfill
\begin{minipage}[t]{0.48\linewidth}
    \textbf{Key Understandings}
    \begin{itemize}
        \item $A$ and $B$ are independent if knowing one does not change the
              probability of the other: $P(A|B) = P(A)$.
        \item If and only if independent: $P(A \cap B) = P(A) \cdot P(B)$.
        \item The union $A \cup B$ means ``$A$ \emph{or} $B$ (or both).''
        \item The addition rule subtracts $P(A \cap B)$ to avoid double-counting.
        \item Mutually exclusive events with positive probability are \emph{dependent},
              not independent.
    \end{itemize}
\end{minipage}
\end{skillbox}

\begin{skillbox}[{Vocabulary, Concepts, \& Theorems}]{greenbox}
\begin{minipage}[t]{0.48\linewidth}
    \begin{itemize}
        \item \textbf{Independent events} --- $P(A|B) = P(A)$; knowing $B$ does not change $P(A)$
        \item \textbf{Dependent events} --- $P(A|B) \neq P(A)$; knowing $B$ changes $P(A)$
        \item \textbf{Multiplication rule for independent events} --- $P(A \cap B) = P(A) \cdot P(B)$
    \end{itemize}
\end{minipage}
\hfill
\begin{minipage}[t]{0.48\linewidth}
    \begin{itemize}
        \item \textbf{Union} $A \cup B$ --- outcomes in $A$ \emph{or} $B$ (or both)
        \item \textbf{Addition rule} --- $P(A \cup B) = P(A) + P(B) - P(A \cap B)$
        \item \textbf{Mutually exclusive $\neq$ independent} --- a critical distinction
    \end{itemize}
\end{minipage}
\end{skillbox}

\begin{skillbox}[Activate Prior Knowledge \& Spiral Review]{sky}
\begin{minipage}[t]{0.48\linewidth}
    \textbf{Prior Knowledge Check (3--4 min)}
    \begin{itemize}
        \item Review: conditional probability formula $P(A|B)$ (Lesson 4.5) and
              mutually exclusive condition $P(A \cap B) = 0$ (Lesson 4.4).
        \item From the diet/health table: ``Is the probability of a healthy outcome
              the same for plant-based and omnivore dieters?'' (foreshadow independence)
        \item Bridge: ``When the answer is YES, we have a special relationship
              called independence.''
    \end{itemize}
\end{minipage}
\hfill
\begin{minipage}[t]{0.48\linewidth}
    \textbf{Connections Forward}
    \begin{itemize}
        \item Independence is essential for probability calculations in Unit 5
              (sampling distributions assume independence or near-independence).
        \item Addition rule extends to complement rule applications.
        \item Union probability appears throughout probability calculations on the AP exam.
    \end{itemize}
\end{minipage}
\end{skillbox}

\begin{skillbox}[Hook \& Lesson]{sky}
\begin{minipage}[t]{0.48\linewidth}
    \textbf{Hook (4--5 min)}\\
    Present these two claims:
    \begin{enumerate}
        \item ``If I flip a fair coin twice, knowing the first flip was heads tells me
              nothing about the second flip.''
        \item ``If I draw two cards from a deck without replacement, knowing the first
              was a heart changes the probability the second is a heart.''
    \end{enumerate}
    Ask: ``What's different about these two situations?''\\[4pt]
    Transition: \textit{``When knowledge of one event doesn't change the probability
    of another, the events are independent.''}
\end{minipage}
\hfill
\begin{minipage}[t]{0.48\linewidth}
    \textbf{Lesson Flow (15--18 min)}
    \begin{enumerate}
        \item Define independence via $P(A|B) = P(A)$; develop equivalent form
              $P(A \cap B) = P(A) \cdot P(B)$.
        \item Check independence from a two-way table: compare $P(A|B)$ to $P(A)$.
        \item Address the big confusion: mutually exclusive (with positive probabilities)
              $\Rightarrow$ dependent, \emph{not} independent.
        \item Define union $A \cup B$; develop the addition rule from the
              double-counting argument.
        \item Work a complete union example from the table.
    \end{enumerate}
\end{minipage}
\end{skillbox}

\begin{skillbox}[Explicit Instruction \& Active Monitoring]{sky}
\begin{minipage}[t]{0.48\linewidth}
    \textbf{Direct Instruction (10 min)}
    \begin{itemize}
        \item Independence check from a table: compute $P(A)$ from the marginal,
              then $P(A|B)$ from the cell; compare.
        \item Show algebraic equivalence: $P(A|B) = P(A) \Leftrightarrow
              P(A \cap B) = P(A) \cdot P(B)$.
        \item Demonstrate the double-counting problem: add $P(A)$ + $P(B)$, then
              subtract $P(A \cap B)$.
        \item Emphasize: for mutually exclusive events, $P(A \cup B) = P(A) + P(B)$
              (a special case, since $P(A \cap B) = 0$).
    \end{itemize}
\end{minipage}
\hfill
\begin{minipage}[t]{0.48\linewidth}
    \textbf{Active Monitoring}
    \begin{itemize}
        \item Watch for the ``independent = mutually exclusive'' confusion.
              Pose directly: ``Can mutually exclusive events with positive probabilities
              be independent? Why not?''
        \item Check addition rule setup: students must subtract $P(A \cap B)$,
              not add it.
        \item Cold-call: ``If $P(A) = 0.3$ and $P(B) = 0.4$ and $A$, $B$ are
              independent, what is $P(A \cap B)$?''
    \end{itemize}
\end{minipage}
\end{skillbox}

\begin{skillbox}[Group Work \& Differentiation]{redbox}
\begin{minipage}[t]{0.48\linewidth}
    \textbf{Group Activity (10--12 min)}\\
    Three-tier activity using a school program $\times$ satisfaction survey table.
    \begin{itemize}
        \item Tier R: identify $P(A)$, $P(A|B)$, and decide if they appear equal.
        \item Tier A: formal independence check; apply the addition rule.
        \item Tier E: compare independence and mutually exclusive; create a
              scenario that demonstrates the distinction.
    \end{itemize}
\end{minipage}
\hfill
\begin{minipage}[t]{0.48\linewidth}
    \textbf{Differentiation}
    \begin{itemize}
        \item \textit{Support}: Provide a fill-in template: ``$P(A) = \_\_\_$,
              $P(A|B) = \_\_\_$. Since these are [equal / not equal], events $A$
              and $B$ are [independent / dependent].''
        \item \textit{Extension}: Prove algebraically that if $A$ and $B$ are
              mutually exclusive and $P(A) > 0$, $P(B) > 0$, then they cannot be
              independent.
        \item \textit{ELL}: Venn diagram showing ``or means everything in either
              circle, subtract the overlap.''
    \end{itemize}
\end{minipage}
\end{skillbox}

\begin{skillbox}[Individual Work \& Assessment]{redbox}
\begin{minipage}[t]{0.48\linewidth}
    \textbf{Exit Ticket (5 min)}\\
    From a two-way table (school preferences survey):
    \begin{enumerate}
        \item Test whether two events are independent.
        \item Compute $P(A \cup B)$ using the addition rule.
        \item Identify whether two specific events are mutually exclusive
              or independent (or neither, or both).
    \end{enumerate}
\end{minipage}
\hfill
\begin{minipage}[t]{0.48\linewidth}
    \textbf{AP-Style Check}\\
    Multiple-choice: ``Events $A$ and $B$ are mutually exclusive. Which of the
    following must be true?''\\
    Options include: independent, $P(A \cap B) = 0$, $P(A \cup B) = P(A) \cdot P(B)$, etc.\\[4pt]
    Answer: $P(A \cap B) = 0$.
\end{minipage}
\end{skillbox}

\begin{skillbox}[Reinforcement \& Extension]{goldbox}
\begin{minipage}[t]{0.48\linewidth}
    \textbf{Homework / Practice}
    \begin{itemize}
        \item Five to six problems using a school preferences two-way table:
              independence tests, addition rule calculations, and the
              mutually exclusive vs.\ independent distinction.
        \item One free-response item requiring complete justification.
    \end{itemize}
\end{minipage}
\hfill
\begin{minipage}[t]{0.48\linewidth}
    \textbf{Extension / Preview}
    \begin{itemize}
        \item \textit{Extension}: Prove that if $P(A) > 0$ and $P(B) > 0$,
              mutually exclusive events cannot be independent.
        \item \textit{Preview}: Lessons 4.7--4.9 apply these rules to random
              variables and probability distributions --- the same rules, with
              variables rather than specific outcomes.
    \end{itemize}
\end{minipage}
\end{skillbox}

\end{document}
